\documentclass[10pt,twocolumn]{article}
\usepackage[letterpaper,margin=0.72in,columnsep=0.22in]{geometry}
\usepackage{amsmath,amssymb,mathtools}
\usepackage{booktabs,tabularx,array}
\usepackage{enumitem}
\usepackage{microtype}
\usepackage[hidelinks]{hyperref}
\usepackage{url}
\usepackage{enumitem}
\usepackage{graphicx}
\usepackage{caption}
\usepackage{fancyhdr}
\usepackage{titlesec}
\usepackage{newtxtext,newtxmath}
\setlength{\parindent}{1em}
\setlength{\parskip}{0pt}
\setlength{\emergencystretch}{2em}
\setlist{nosep,leftmargin=*}
\titleformat{\section}{\bfseries\large}{\thesection}{0.5em}{}
\titleformat{\subsection}{\bfseries\normalsize}{\thesubsection}{0.5em}{}
\titleformat{\subsubsection}{\bfseries\normalsize}{\thesubsubsection}{0.5em}{}
\pagestyle{fancy}
\fancyhf{}
\fancyhead[LE,RO]{RAI: Value Transfer with Account-Local Consensus}
\fancyfoot[C]{\thepage}
\newcommand{\NC}{\mathrm{NC}}
\newcommand{\FC}{\mathrm{FC}}
\newcommand{\FF}{\mathrm{FF}}
\newcommand{\TC}{\mathrm{TC}}
\newcommand{\ET}{\mathrm{ET}}
\newcommand{\Hash}{\mathrm{H}}
\newcommand{\StateHash}{\mathrm{H}_{\mathrm{state}}}
\newcommand{\VoteHash}{\mathrm{H}_{\mathrm{vote}}}
\newcommand{\DeriveCommittee}{\mathrm{DeriveCommittee}}
\newcommand{\First}{\mathrm{First}}
\newcommand{\Notar}{\mathrm{Notar}}
\newcommand{\Final}{\mathrm{Final}}
\newcommand{\Irrev}{\mathrm{Irrev}}
\newcommand{\Commit}{\mathrm{Committed}}
\newcommand{\FinalCert}{\mathrm{FinalCert}}
\newcommand{\Include}{\mathrm{include}}
\newcommand{\ReconReq}{\mathrm{ReconReq}}
\newcommand{\ReconReply}{\mathrm{ReconReply}}
\newcommand{\BuildState}{\mathrm{BuildState}}

\title{\vspace{-1.1cm}\textbf{RAI: Value Transfer with Account-Local Consensus}}
\author{Anonymous Author(s)}
\date{}

\begin{document}
\maketitle
\vspace{-0.7em}

\begin{abstract}
We study value transfer between single-owner accounts under a changing committee of Byzantine validators. RAI runs one certified block tree per account and adapts Kudzu's first, notarization, final, and second-look voting rules to account-local elections. Account elections have no slot timer and no timeout value: an owner may broadcast a valid block at any time while an epoch is open, and validators vote whenever the block becomes available. An account block may contain multiple sends and receives. Its hash contains neither an epoch nor a voting-attempt identifier, and its account slot is fixed by its parent.

RAI separates notarized ledger membership from application effects. Only finalized blocks affect balances, received-send sets, or committee derivation. At an epoch boundary, a decided epoch state can finalize a uniquely included block, checkpoint-notarize multiple conflicting included blocks without applying them, or omit provisional blocks and thereby roll them back. If no block survives at an unresolved slot, the next epoch begins a clean election for that slot. If conflicting notarized blocks survive at slot $v$, recovery advances the account chain: after the earlier epoch is closed, the owner creates a new block at slot $v+1$ that names one surviving block as parent. Finalization of the child selects that branch and recursively finalizes its unresolved notarized ancestors.

Finalization therefore has an ancestral closure rule. A finalized block in epoch $e$ may finalize unresolved notarized ancestors from the same epoch immediately. If an unresolved ancestor was voted in an earlier epoch $e'\ne e$, epoch $e'$ must first be closed and that ancestor must survive as checkpoint-notarized state. This preserves intra-epoch pipelining while making epoch closure the recovery barrier for equivocated accounts.

At an epoch boundary, validators stop issuing old-epoch account votes and sign compact reports. Reports commit to certified account state and to the reporter's own residual signed votes not yet summarized by that state. Reconciliation is state based: a requester names a locally available source root and a historical target root, and a replica that knows both returns a reconstructive authenticated difference containing the state edits and hash material needed to recover the target. The requester accepts the reconstruction only if recomputing the certified-state hash yields the signed target root; there is no separate reconciliation-proof object. A joint epoch election under the old and new committees decides the checkpoint ledger and its finalized application projection. The epoch proposal carries selected report roots and only a hash of the deterministically derived state. We give safety arguments and separate progress conditions for account updates, reconciliation, and epoch closure. Implementation and experimental results are reserved for an RsNano-based comparison; no performance results are claimed in this draft.
\end{abstract}

\noindent\textbf{Draft status.} Protocol and proof draft. Sections~4 and~5 reserve space for missing implementation, evaluation, and related-work material.

\section{Introduction}
A transfer system must prevent an account from spending the same value twice while allowing unrelated accounts to progress independently. Reconfiguration adds a second problem: a new validator committee must preserve certified account state accumulated under the old committee without imposing one global transaction order or discarding unresolved forks.

RAI uses one chain position, or \emph{account slot}, per account update. Only the account owner can authorize a block for that slot. A block may batch several sends and receives, and independent accounts need not share a total order. Validators use Byzantine voting to certify competing blocks locally and use a separate epoch-election protocol to checkpoint compatible branches and derive future committees.

Our goals are safety, account progress, low transaction latency, high throughput, and resilience to faulty validators and committee changes. Safety means that finalized account histories never conflict, finalized account decisions remain irreversible across later checkpoints, balances remain nonnegative, and one send cannot be received twice. Progress is stated for owners that eventually stop equivocating and continue gossiping a valid branch. There is no account scheduler or assigned proposal opportunity: an owner may broadcast its block whenever the relevant epoch is open.

\paragraph{Design.} RAI adapts Kudzu~\cite{Kudzu} in four steps.

First, each account slot is a separate voting domain. The owner can broadcast a valid block at any time, and that block itself is the account-election proposal. Account elections have no slot timer, timeout vote, or timeout-generated empty account block. Correct validators simply stop issuing account votes when the epoch reaches its boundary. If an unresolved slot has no block in the decided epoch state, a later epoch starts a clean election at that same slot.

Second, the client block is independent of the voting attempt. A block fixes the account, parent, and an ordered batch of operations. Its hash does not contain the epoch. Votes contain the epoch and slot, so old election votes cannot be replayed as fresh evidence. A block omitted from a decided checkpoint may be proposed again in a later clean election, but conflicting checkpoint-notarized blocks are recovered differently: the owner advances to the next slot with a new child that explicitly chooses one surviving parent.

Third, RAI separates notarized ledger membership from finalization and application effects. A decided epoch state contains a checkpoint ledger that may include conflicting notarized blocks. If exactly one compatible block survives at a slot, checkpoint inclusion finalizes it; if multiple incompatible blocks survive, they are checkpoint-notarized and remain provisional. Only finalized blocks affect balances, receive-once state, or committee derivation. A finalization certificate for a descendant recursively finalizes the unresolved notarized ancestor prefix on the same branch. Same-epoch ancestors can be finalized immediately by such a descendant; an ancestor from an earlier epoch must first belong to a closed epoch state. Thus epoch closure is the barrier that makes cross-epoch fork recovery safe.

Fourth, epoch reports are state commitments rather than complete vote logs. Each reporter signs a certified-state root and a residual-vote root committing only to that reporter's own signed account votes that are not yet summarized by the corresponding certified status. Votes already represented by a verifiable notarization or finalization status need not be duplicated in the residual commitment. Replicas continue adding newly learned signed votes after reporting and update vote-notarized or vote-finalized status whenever the corresponding certificate becomes locally constructible. Reconciliation uses common descendant states: a request names a locally available source root and a historical target root, and a replica that knows both supplies the reconstructive difference. The receiver verifies reconciliation by reconstructing the target and recomputing its signed root, not by checking a separate reconciliation proof. No replica is required to recursively compare two report dictionaries.

The epoch election is one logical election run by the old and new committees. Unlike account elections, the epoch-election Kudzu instances retain their voting timeouts. Matching locally constructed notarization and finalization certificates are required for success. A timeout certificate in either committee abandons the epoch-election slot. If the two committees notarize different epoch values, those conflicting locally constructed certificates are also safe evidence to skip that slot.

\paragraph{Scope and contributions.} RAI combines timer-free account-local voting, batched account operations, checkpoint-notarized fork state, ancestral finalization, child-based recovery from equivocation, state-root reporting, convergence-based historical-state reconciliation, local epoch-preservation locks, and lagged joint reconfiguration. We give safety arguments and distinguish account progress, report availability, and epoch-election progress rather than hiding them behind one liveness statement. The intended prototype and comparison use RsNano~\cite{RsNano}; Section~4 reserves space for implementation and measurements.

\section{Preliminaries}
\subsection{System model}
Clients own accounts and authorize blocks. Validators check blocks and vote. Every committee has
\begin{equation}
N = 3f + 2p + 1
\end{equation}
validators, of which at most $f$ are Byzantine. The parameter $p\ge 0$ controls the fast-path threshold. Corruptions are static, and the fault bound holds separately for every derived committee. Votes have equal weight. Offline validators consume the unavailable budget assumed by a liveness statement; they are not extra Byzantine faults.

Validator identities have known public keys. Signatures are unforgeable and hashes are collision resistant. RAI uses individually verifiable signatures. A certificate is a local quorum view: a validator constructs it from a set of distinct signed votes satisfying the relevant quorum predicate. Certificates are not separately signed or gossiped as protocol objects; validators gossip the underlying signed votes and independently construct the same certificate whenever they collect sufficient support.

\paragraph{Membership.} A closed epoch state is $S_e=(L_e,\Sigma_e)$. $L_e$ is the authenticated checkpoint ledger: an account forest containing finalized blocks and any unresolved checkpoint-notarized forks that the epoch decision preserves. $\Sigma_e$ is the deterministic application projection of $L_e$ obtained by replaying finalized blocks only; it records finalized account histories, balances, and received-send identifiers. Provisional notarized forks in $L_e$ have no application effect and do not influence committee selection. The committee derived from epoch $e$ is
\begin{equation}
C_e = \DeriveCommittee(\Sigma_e).
\end{equation}
Derivation is deterministic and reads no provisional fork state. Genesis defines the initial checkpoint ledger $L_G$, application state $\Sigma_G$, $S_G=(L_G,\Sigma_G)$, and $C_G$. Set $S_{-1}=S_G$ and $C_{-2}=C_{-1}=C_G$. Committee admission and the mechanism enforcing the per-committee fault bound are deployment assumptions.

\paragraph{Gossip and availability.} All transmitted protocol objects share one peer-to-peer gossip layer, in the style of Tendermint~\cite{Tendermint}: full account blocks, signed votes, report roots and their referenced evidence, state-difference requests and replies, epoch proposals, and lock-objection evidence. Certificates are not transmitted; each validator constructs them locally from the signed votes it has received. Any valid object sent or received by a correct validator eventually reaches all correct validators. During a $\delta$-synchronous interval, propagation takes at most $\delta$ real time.

Correct validators retain and serve the data needed to verify certified blocks, reconstruct historical report roots, recover residual vote evidence, check local lock objections, and apply state differences. A correct validator has the full block before first-voting for it. Missing data are fetched; absence of a reply is not proof of invalidity. Correct reporters retain enough persistent state history or authenticated deltas to reconstruct every certified-state root they sign, including the hash material needed to recompute that root from an available source state.

\paragraph{Time.} Safety does not require a message-delay bound. Liveness requires recurring sufficiently long synchronous intervals, following partial synchrony~\cite{DLS,Kudzu}. Correct clocks need not agree on readings. During a real interval of length $t$, each advances by between $(1-\rho)t$ and $(1+\rho)t$, for $0\le \rho<1$. Account elections have no local slot timeout. The timeout $\Delta_E$ is used only by the joint epoch-election Kudzu instances and satisfies
\begin{equation}
\frac{\Delta_E}{1+\rho} \ge 2\delta.
\end{equation}
An open account epoch lasts $D$ local time. If correct epoch openings differ by at most $\kappa$ real time, their common-open interval has length at least $D/(1+\rho)-\kappa$. Account-liveness statements require a sufficiently long common-open portion for block and vote propagation before correct validators stop issuing that epoch's account votes.

\subsection{State model}
The ledger contains genesis and one account forest per account. Every block has an account slot determined by its parent depth: if the parent is at slot $v-1$, the block is at slot $v$. Slots are account positions, not voting attempts. Epoch closure never creates an empty block. If a decided epoch state contains no surviving block at an unresolved slot, the next epoch begins a clean election for that same slot. If the decided state preserves conflicting blocks at slot $v$, those blocks occupy slot $v$ provisionally and recovery may proceed at slot $v+1$ by extending one of them.

An account block has the form
\begin{equation}
B=(a,h_p,\mathbf{x},\sigma_a),\qquad h(B)=\Hash(\mathrm{Block},a,h_p,\mathbf{x}),
\end{equation}
where $h_p$ is the parent account-block hash and $\mathbf{x}=(x_1,\ldots,x_m)$ is an ordered finite batch of operations. Neither an epoch nor a voting-attempt identifier appears in $h(B)$. The owner signature $\sigma_a$ authenticates the block contents.

For the transfer interface, each $x_j$ is either $\mathrm{Send}(b,z)$, which sends positive amount $z$ to account $b$, or $\mathrm{Receive}(s)$, which receives one prior send. A send in position $j$ of block $B$ has identifier
\begin{equation}
s=(h(B),j).
\end{equation}
Operations are replayed in block order. A send decreases the sender's running balance. A receive increases the destination balance by the amount named by $s$. A valid block has a valid owner signature and parent, never makes the running balance negative when its branch is finalized, and does not receive the same send twice either in the finalized parent history or earlier in the same block. Every receive names a send addressed to this account and carries the finalized-send evidence required by the protocol.

Two distinct valid blocks at the same account slot are \emph{conflicting}: both occupy one account position, so neither can lie on the other's ancestry. Direct siblings with the same parent are the simplest case, and blocks whose parents already conflict are the general one. An owner equivocates when it authenticates incompatible blocks for the same account position or incompatible continuations of an unresolved fork. An owner is non-equivocating after time $T$ if every block it authenticates after $T$ lies on one parent chain.

\paragraph{Vote-notarized, checkpoint-notarized, finalized, and irreversible state.} A block is \emph{vote-notarized} in domain $(e,a,v)$ when a validator can locally construct $\NC(h(B))$ from epoch-$e$ signed account votes. A block is \emph{checkpoint-notarized} when a decided epoch state includes it at a slot together with one or more conflicting included blocks. Either status is called \emph{notarized}. Notarization alone has no application effect.

A block is \emph{finalized} in either of two ways. First, a valid account-domain fast or normal finalization certificate may finalize the block and its eligible unresolved ancestor prefix. Second, a decided epoch state checkpoint-finalizes a uniquely surviving block and the same eligible ancestor prefix. Finalized blocks are irreversible. $\Sigma_e$ contains only finalized application state and therefore never rolls back, while provisional checkpoint-notarized branches in $L_e$ may later be removed.

\paragraph{Ancestral finalization closure.} Let $B$ be a complete block voted in epoch $e$. Starting from $B$ and walking backward, form the maximal contiguous prefix of unresolved notarized ancestors until the first already-finalized ancestor. An unresolved ancestor voted in the same epoch $e$ may belong to this prefix immediately. If an unresolved ancestor was voted in an earlier epoch $e'\ne e$, then epoch $e'$ must already be closed and the ancestor must survive in that decided state as checkpoint-notarized. Every external receive dependency of every block in the prefix must already be finalized. When $B$ becomes finalized, every block in this eligible unresolved prefix becomes finalized atomically in parent order. Thus a descendant may select one branch of an earlier closed fork, but it cannot use its own finalization to reach backward through an earlier epoch that has not yet closed.

A valid epoch state may never remove a finalized block or preserve a block conflicting with a finalized account position. When finalization of a descendant selects one branch, incompatible provisional siblings and their provisional descendants are rolled back. Only finalized blocks contribute operations to balances or received-send sets.

\section{RAI}
\subsection{Protocol}
\subsubsection{Account elections}
For an account $a$, account slot $v$, and epoch $e$, account votes belong to domain $(e,a,v)$. There is no separate proposal message and no account-slot timer. The owner gossips the signed block itself, and validators vote on its hash while epoch $e$ is open. The slot is derived from the parent; the epoch is supplied only by the voting domain. If an unresolved slot is later clean because the decided epoch state included no block there, the owner may submit the same valid block or a different valid block at that slot in a later epoch, but every new epoch uses fresh votes.

There are first, notarization, and final votes. A first vote is one signed message. A signed $\First(h)$ counts both as that validator's first vote and as one notarization contribution for $h$. After the second-look rule, a validator may also issue $\Notar(h)$. When constructing a quorum, a validator identity counts at most once for a value even if both a first vote and a later notarization vote from that identity are known. The thresholds are
\begin{equation}
q=N-f-p=2f+p+1,\qquad q_{\mathrm{fast}}=N-p.
\end{equation}
Within one account voting domain, $\NC(h)$, $\FC(h)$, and $\FF(h)$ denote locally constructed vote-notarization, normal-finalization, and fast-finalization certificates. A normal notarization certificate $\NC(h)$ consists of $q$ distinct validators supporting $h$, where each support may be either $\First(h)$ or $\Notar(h)$. A fast certificate uses $q_{\mathrm{fast}}$ distinct $\First(h)$ votes. Normal finalization uses $q$ distinct $\Final(h)$ votes.

\begin{table}[t]
\centering
\caption{Quorums. The example is a parameter calculation, not an experiment.}
\begin{tabular}{lcc}
\toprule
Predicate & Threshold & $f=p=1$\\
\midrule
Account notarization & $q$ & 4\\
Normal account finalization & $q$ & 4\\
Fast account finalization & $q_{\mathrm{fast}}$ & 5\\
Selected reports & $N-f$ & 5\\
Residual-vote checkpoint inclusion & $1$ & 1\\
\bottomrule
\end{tabular}
\end{table}

\paragraph{Parents and slot progression.} There are no skipped account slots. A block at slot $v>1$ names a concrete parent at slot $v-1$. A notarized parent may provisionally support a child at the next slot. If epoch $e$ closes with conflicting checkpoint-notarized blocks $A_v$ and $B_v$, recovery does not retry either block at slot $v$. Instead, after that epoch is closed, the owner may create a new block at slot $v+1$ naming one survivor, for example $A_v$, as parent. Finalization of that child applies the ancestral finalization closure, finalizing $A_v$ and eliminating $B_v$ and its incompatible provisional descendants. If no block survives at slot $v$, the next election remains a clean election for slot $v$.

A placement $(e,a,v,h(B))$ is \emph{complete} when the validator has the full valid block, valid parent and dependency evidence, and can locally construct $\NC(h(B))$ from signed votes in domain $(e,a,v)$. Completeness adds vote-notarized status to the certified tree. It does not apply the block to balances.

\paragraph{Finalization eligibility.} A complete block $B$ voted in epoch $e$ is finalization-eligible when every block in its ancestral finalization closure is valid and notarized, every external receive dependency of that closure is finalized, and every unresolved ancestor from an earlier epoch belongs to an already closed epoch state. Unresolved notarized ancestors from epoch $e$ itself need not wait for the epoch checkpoint. A locally constructed $\FF(B)$ or $\FC(B)$ finalizes $B$ and the entire eligible unresolved ancestor prefix. If an earlier-epoch ancestor has not yet closed, $B$ may still be notarized, but correct validators do not final-vote for it yet.

\paragraph{Voting loop.} Algorithm~1 states the timer-free account rules. Each validator sends at most one first vote and one final vote per domain, and final-votes $h$ only while it has issued notarization support for no block other than $h$. Let $\mathit{many}$ be the hashes with at least $f+p+1$ recorded first votes. After first-voting, a validator takes a second look at each new hash in $\mathit{many}$, fetches and validates the block, and may issue additional notarization support. There is no account timeout value and no split-vote timeout rule; unresolved competition is carried to the epoch decision if it is not resolved by finalization first.

Data fetches do not block other enabled domains. After finalizing and locally exiting a domain, a correct validator continues receiving and relaying old blocks and signed votes but issues no new vote in that domain. At the epoch boundary it permanently stops issuing all epoch-$e$ account votes and signs its report. Inactive accounts require no timer and no scheduled service.

\begin{figure*}[t]
\small
\centering
\fbox{\begin{minipage}{0.94\textwidth}
\textbf{Algorithm 1: One timer-free account voting domain $(e,a,v)$ at a validator}\\[0.3em]
1. Keep the domain dormant until a valid owner block for slot $v$ or a valid vote for $(e,a,v)$ is observed.\\
2. Gossip valid blocks and signed votes; never gossip a certificate as a separate protocol object.\\
3. Record the first vote received from each sender.\\
4. If not first-voted and a valid owner block $B$ for slot $v$ is available, send $\First(h(B))$; this same signed vote also counts toward $\NC(h(B))$.\\
5. For each unchecked $h\in\mathit{many}$, fetch and validate the block; if valid and additional support is useful, send $\Notar(h)$.\\
6. If a placement $h$ is complete, finalization-eligible, and this validator has issued notarization support for no block other than $h$, send $\Final(h)$. If $\FF(h)$ or $\FC(h)$ becomes locally constructible, finalize $h$ and its eligible unresolved ancestor prefix and exit the domain.\\
7. Otherwise continue gossip and voting until finalization or the epoch boundary. At the boundary, stop issuing epoch-$e$ account votes; unresolved blocks are handled by the decided epoch state.
\end{minipage}}
\caption{Account-domain voting has no slot timer or timeout value. First votes are notarization support; unresolved forks are resolved or preserved by the epoch decision.}
\end{figure*}

\subsubsection{Epochs and lagged committees}
An epoch is open, closing, or closed. While open, its committee may issue account votes in any activated account domain. At the end of local duration $D$, each correct validator permanently stops issuing epoch-$e$ account votes and signs its report. A joint epoch-election finalization certificate, constructed locally from votes in both epoch-election committee domains, later makes epoch $e$ closed and fixes $S_e=(L_e,\Sigma_e)$.

The committee running epoch $e$ is
\begin{equation}
K_e=C_{e-2}.
\end{equation}
The epoch election for $e$ uses old committee $O_e=C_{e-2}$ and new committee $N_e=C_{e-1}$. Its decided finalized application state $\Sigma_e$ determines $C_e$, which first runs account elections in $e+2$. For $e=0$, both epoch-election committees are $C_G$, but their voting domains remain distinct.

Epoch $e+1$ need not wait for $S_e$ to begin receiving and vote-notarizing account blocks because it uses $C_{e-1}$. It may also receive a child that references an unresolved block voted in epoch $e$. However, that child cannot become finalization-eligible through the earlier parent until epoch $e$ is closed and the parent survives the decided checkpoint ledger as checkpoint-notarized or finalized. Thus account notarization may overlap epoch closure, while cross-epoch finalization observes a strict closure barrier.

Within one epoch, finalization may pipeline through several unresolved notarized account slots: finalization of a descendant finalizes its eligible same-epoch ancestor prefix without checkpoint delay. Across epochs, the earlier epoch must close first. In particular, if epoch $e$ closes with conflicting checkpoint-notarized blocks $A_v$ and $B_v$, a later recovery block $A_{v+1}$ may name $A_v$ as parent; finalization of $A_{v+1}$ then finalizes $A_v$ and rolls back $B_v$. If the decided epoch state contains no block at an unresolved slot, later work resumes with a clean election at that slot.

\begin{figure}[t]
\centering
\small
\begin{tabular}{c|c|c}
Account elections & Committee & Epoch-election relation\\
\hline
epoch $e$ & $K_e=C_{e-2}$ & election $e$: $C_{e-2}$ and $C_{e-1}$\\
epoch $e+1$ & $K_{e+1}=C_{e-1}$ & election $e+1$: $C_{e-1}$ and $C_e$\\
epoch $e+2$ & $K_{e+2}=C_e$ & $S_e$ derives $C_e$\\
\end{tabular}
\caption{Lagged reconfiguration. Account notarization may overlap a predecessor epoch election, but finalization across an epoch boundary waits for the earlier epoch to close.}
\end{figure}

\subsubsection{Certified-state reports and reconciliation}
At its epoch boundary, each correct old-committee validator signs one report after stopping epoch-$e$ account voting. The report does not contain a complete per-vote dictionary. Instead it commits to two authenticated objects:
\begin{enumerate}
\item a certified-state root $r_i=\StateHash(T_i)$ for the reporter's canonical epoch-$e$ certified block tree, including every complete notarized block known to the reporter and any fast or normal finalization status already known for those blocks; and
\item a residual-vote root $g_i=\VoteHash(G_i)$ for reporter-local vote evidence not yet reflected by the corresponding certified status in $T_i$.
\end{enumerate}
The signed report binds the session, epoch, old-committee identifier, reporter identity, $r_i$, and $g_i$.

The residual vote evidence has a narrow role even when its concrete representation is large. $G_i$ is not a log of every vote the reporter has received. It contains only reporter $i$'s own signed account votes whose effect is not yet summarized by the corresponding certified status in $T_i$. Thus, for a block with no $\NC$ entry in $T_i$, $G_i$ contains reporter $i$'s signed $\First(h)$ if it issued one and any later $\Notar(h)$ that is still unsummarized. A first vote is itself notarization support, but first votes from other validators are not copied into $G_i$ merely because reporter $i$ received them. For a block that is represented as notarized but not finalized in $T_i$, $G_i$ contains reporter $i$'s own unsummarized $\Final(h)$, if any. Once a verifiable $\NC$, $\FF$, or $\FC$ status is represented in the certified state, the individual votes summarized by that status need not also be duplicated in a later residual view. The signed report's original $G_i$ remains immutable and reconstructible. These residual records preserve reporter-local signed support from which certificates may become locally constructible only after several reports are combined.

\paragraph{Local epoch-preservation locks.} In addition to signing its report, a correct old-committee validator retains its own epoch-$e$ account-vote history until the epoch election for $e$ decides. When validating an epoch proposal and its locally derived checkpoint ledger, it checks that every block hash for which it issued an account vote is either included in $L_e$ or excluded by verifiable incompatibility with a finalized branch. Inclusion need not finalize the block: if conflicting included blocks survive at that slot, the epoch decision checkpoint-notarizes them without application effects. If this local preservation check fails, the validator does not vote for the epoch proposal and gossips the signed evidence that caused the rejection. A later leader can include the objecting validator's report, making that voted block report-visible to the deterministic derivation.

\paragraph{State evolution after reporting.} Signing a report freezes the historical roots $r_i$ and $g_i$, not the validator's live certified tree. Correct replicas continue to receive old epoch-$e$ blocks and signed votes through gossip. When newly received votes make a certificate locally constructible, they add the certified block or stronger finalization status to their live canonical tree and obtain a new state root. They issue no new epoch-$e$ votes after their report.

A correct validator retains enough authenticated history to relate a signed historical root to later roots. For two certified states known by a responder, a state difference $\Delta(r_s,r_t)$ contains the canonical edits needed to reconstruct the state committed by target root $r_t$ from the state committed by source root $r_s$. These edits may add or remove blocks, change certified status, and carry any sibling hashes or other authentication material required by the chosen authenticated representation. This authentication material is part of the difference itself; RAI defines no separate reconciliation-proof object. Implementations may realize $\Delta$ with a persistent Merkle structure, an append-only authenticated log, a versioned trie, or another canonical authenticated representation.

\paragraph{Reconciliation.} To reconstruct target report root $r_t$, a validator that currently has certified state $T_s$ with root $r_s$ gossips
\begin{equation}
\ReconReq(r_s,r_t).
\end{equation}
A replica responds only if it knows both states and can return a reconstructive difference,
\begin{equation}
\ReconReply(r_s,r_t,\Delta(r_s,r_t)).
\end{equation}
The requester applies the difference to $T_s$ to obtain a candidate historical state $T_t$ and accepts that reconstruction exactly when
\begin{equation}
\StateHash(T_t)=r_t.
\end{equation}
Thus successful root reconstruction is the reconciliation check. No additional signature or proof over $\Delta$ is required beyond the signed report commitment to $r_t$ and the authenticated hash material needed to recompute it. If no replica currently knows both source and target states, no response is required.

As gossip equalizes the certified evidence held by correct replicas, their live roots repeatedly acquire common descendant states. If a direct request initially has no responder, the requester may later reissue it using a newer locally available common-descendant root as $r_s$. A replica retaining the historical target and that descendant can then supply $\Delta(r_s,r_t)$, allowing the requester to reconstruct the historical state and recompute $r_t$. If old-epoch evidence eventually quiesces, all correct replicas converge to the same final certified block tree and state root; that final root is a universal source state from which every retained correct report root can be reconstructed. The same gossip layer fetches block bodies and the constituent signed votes needed to verify any $\NC$, $\FF$, or $\FC$ status represented in the reconstructed state, as well as lock-objection evidence. Those certificate votes justify certified status; they are distinct from reconciliation authentication material.

A report is usable only after both committed objects have been recovered and checked by hash recomputation: its certified state $T_i$ must satisfy $\StateHash(T_i)=r_i$, and its residual object $G_i$ must satisfy $\VoteHash(G_i)=g_i$. The reconstructed certified statuses must additionally be supported by the required signed account votes. A signed root by itself does not make unavailable contents usable. A Byzantine reporter may sign an unreconstructible root or withhold the committed contents; such a report simply never becomes usable. The epoch election needs only $N-f$ usable reports, so correct reporters suffice under the availability assumptions.

\subsubsection{Deriving the epoch state from reports}
An epoch proposal selects exactly
\begin{equation}
q_{\mathrm{report}}=N-f=2f+2p+1
\end{equation}
usable reports from distinct old-committee validators. Let $Q_e$ denote the canonically ordered set of their signed report identifiers and roots. The epoch proposal does not carry a serialized candidate state. Instead, every validator reconstructs the reports in $Q_e$ and deterministically derives
\begin{equation}
S_e(Q_e)=\BuildState(S_{e-1},Q_e)=(L_e,\Sigma_e),
\end{equation}
then computes
\begin{equation}
d_e=\StateHash(S_e(Q_e)).
\end{equation}
The state field carried by the epoch proposal is only $d_e$. Thus two validators that accept the same proposal independently reconstruct the selected reports, derive the same checkpoint ledger and finalized application projection, and obtain the proposed hash.

For a block hash $h$, first inspect the reconstructed certified states in $Q_e$. If a selected state records a vote-notarized status and the fetched constituent signed votes satisfy the $\NC$ predicate, then $h$ is \emph{certified-visible}. If a selected state records a fast or normal account-finalization status and the fetched signed votes satisfy the corresponding $\FF$ or $\FC$ predicate, then $h$ is \emph{final-visible}. These certificate checks validate the meaning of the reconstructed state; they are separate from the hash-reconstruction check that binds the state to its report root.

For a block not certified-visible, let $M_Q(h)$ count distinct selected reporters whose reconstructed residual object contains that reporter's own epoch-$e$ first or notarization support for $h$. Votes merely received from other validators do not contribute through that reporter's residual object. Define
\begin{equation}
\Include_Q(h) \Longleftrightarrow h\text{ is certified-visible}\ \lor\ M_Q(h)\ge 1.
\end{equation}
A single selected correct reporter's signed account vote is therefore enough to make the block a report-derived checkpoint candidate. This does not give the block application effects. If it survives together with a conflict, the epoch decision only checkpoint-notarizes it; if it is the unique compatible survivor, the joint epoch decision may checkpoint-finalize it.

The threshold of one is deliberate and has a price. Lemma~3.7 guarantees only $p+1$ correct reporters in the relevant intersection, which is a single reporter when $p=0$, so no larger threshold would preserve it. In exchange, a selected Byzantine reporter can record residual support for any block it chooses and force that block into $L_e$. Rule~2 confines this to blocks the owner itself authenticated, so an owner that never equivocates is unaffected; against an owner that has equivocated, one faulty validator can keep a second valid block alive at the slot, which denies the slot the unique survivor of rule~4 and defers the account to child-based recovery at the next slot. The cost is bounded by the owner's own equivocation and does not reach the finalized projection.

Residual final-vote evidence is also reconciled. A report-visible account-finalization certificate is mandatory for the derivation. Stronger evidence learned outside $Q_e$ does not silently change $S_e(Q_e)$; a correct validator whose local preservation rule detects an omitted voted block rejects that epoch proposal and gossips the signed reason. Selecting the objecting validator's report in a later proposal exposes at least one signed support and therefore makes the block satisfy $\Include_Q$.

The deterministic derivation first forms a candidate account forest from all valid blocks required by predecessor finalized state, final-visible evidence, $\Include_Q$, and the represented ancestry of those blocks. It then computes a unique fixed point under the following rules:
\begin{enumerate}
\item Every block finalized in $\Sigma_{e-1}$ remains finalized and included. No decided epoch state rolls back finalized application state.
\item Every included block has its full body, valid owner authentication, a valid parent chain, and valid external dependencies. A block incompatible with an already finalized account position is excluded.
\item Every final-visible block is finalized together with its eligible unresolved ancestor prefix. This may remove incompatible provisional siblings and their provisional descendants.
\item After mandatory finalization is applied, any included block that is the unique surviving block at its account slot is checkpoint-finalized together with its eligible unresolved ancestor prefix. This rule is repeated to a fixed point, so a unique included descendant may select and finalize a previously unresolved ancestor branch.
\item If two or more incompatible included blocks remain at the same account slot after the finalization fixed point, each surviving block is checkpoint-notarized. These blocks belong to $L_e$ but have no application effect.
\item Every epoch-$e$ block satisfying $\Include_Q$ is included unless it is excluded by incompatibility with a finalized branch. The same condition is enforced for hashes protected by a validating correct validator's local epoch-preservation lock.
\item A provisional block present in an earlier checkpoint ledger but absent from $L_e$ is rolled back, together with any provisional descendant that depends on it. Omission never rolls back a finalized block.
\item $\Sigma_e$ is obtained by replaying finalized blocks only, in account-parent order and operation order. Already finalized operations are not applied twice.
\end{enumerate}

The resulting slot behavior is therefore explicit. If one block survives at a slot, it is finalized. If several conflicts survive, they are checkpoint-notarized without application effects. If no block survives at an unresolved slot, the next epoch begins a clean election at that slot. If several conflicts survive at slot $v$, a later owner block at slot $v+1$ may reference one of them; after the earlier epoch is closed, finalization of that child selects and finalizes the chosen ancestor branch.

\begin{figure}[t]
\centering
\small
\fbox{\begin{minipage}{0.92\columnwidth}
Epoch $e$ may close with conflicting checkpoint-notarized blocks $A_v$ and $B_v$. Neither has application effects. After epoch $e$ is closed, the owner may create $A_{v+1}$ with parent $A_v$. If $A_{v+1}$ finalizes, ancestral closure finalizes $A_v$ and $A_{v+1}$ and rolls back $B_v$ and incompatible provisional descendants.
\end{minipage}}
\caption{Equivocation recovery advances to the next account slot and selects one surviving notarized parent.}
\end{figure}

\subsubsection{The joint epoch election}
An epoch value
\begin{equation}
X=(h_p,Q_e,d_e)
\end{equation}
names a parent epoch value, the selected report roots, and only the hash $d_e$ of the state deterministically derived from those reports. The full $S_e(Q_e)$ is not carried in the epoch proposal. A deterministic epoch leader for epoch-election slot $w$ signs an epoch proposal binding $(e,w,h(X))$. Epoch-election slots are distinct from account slots.

The epoch election has one logical tree and two domain-separated Kudzu voting instances, one under $O_e$ and one under $N_e$. These epoch-election instances retain Kudzu's timeout value and local timeout $\Delta_E$; this timeout mechanism is not used by account elections. A validator present in both committees votes separately in both epoch-election domains. In each domain, a first vote counts toward notarization, timeout votes count toward $\TC$, and certificates are constructed locally from signed votes.

Before voting for $X$, a validator reconstructs every report in $Q_e$, computes $S_e(Q_e)=\BuildState(S_{e-1},Q_e)$, checks $\StateHash(S_e(Q_e))=d_e$, and applies its local epoch-preservation check. Both committees therefore validate the same report-derived epoch value. For an epoch-value hash $h$, define
\begin{align}
\mathrm{JointNotar}(h)&=\NC_O(h)\wedge\NC_N(h),\\
\mathrm{JointFast}(h)&=\FF_O(h)\wedge\FF_N(h),\\
\mathrm{JointFinal}(h)&=\FC_O(h)\wedge\FC_N(h).
\end{align}
Each component certificate is constructed locally from signed votes in its committee domain. A valid epoch placement is complete only with $\mathrm{JointNotar}$. It decides $S_e(Q_e)$ only when complete and accompanied by $\mathrm{JointFast}$ or $\mathrm{JointFinal}$. A certificate in only one committee does not decide the epoch.

An epoch-election slot is abandoned when either committee times out or the two committees notarize different hashes. Formally,
\begin{align}
\ET(e,w) \Longleftrightarrow {}& \TC_O(e,w)\lor\TC_N(e,w)\\
&\lor\exists h\ne h':\NC_O(e,w,h)\wedge\NC_N(e,w,h').
\end{align}
The last clause is the cross-committee conflict timeout. Once a validator has locally constructed the two conflicting notarization certificates from signed votes, they are sufficient skip evidence. A slot can satisfy $\mathrm{JointNotar}$ for one value and $\ET$ at the same time, when a domain has certified a second value as well; Lemma~3.10 shows that such a slot can no longer be decided, so abandoning it discards nothing and the two conditions never disagree about an outcome.

A later epoch proposal must provide $\ET$ evidence for every epoch-election slot between its cited joint-complete parent and itself. A descendant of a non-timeout epoch value carries the same $(Q_e,d_e)$ as its parent; only the lineage hash changes. Hence a completed epoch-election lineage cannot change the selected reports or the report-derived epoch-state hash.

\subsection{Correctness}
The account-domain arguments adapt Kudzu's quorum counting while removing account-domain timeout votes. Account voting is timer-free; unresolved forks are handled by the epoch decision. RAI additionally uses ancestral finalization, checkpoint-notarized fork state, and a local epoch-preservation rule. Safety does not depend on synchrony. Liveness additionally uses the timing and availability assumptions stated below, but no account-service schedule.

\subsubsection{Account-domain safety}
\paragraph{Lemma 3.1 (Finalization excludes conflicting vote-notarization).} If a valid block $h$ is account-finalized in domain $(e,a,v)$, no different block can obtain a vote-notarization certificate in that domain.

\emph{Proof.} For normal finalization, suppose $\FC(h)$ and $\NC(h')$ both exist with $h'\ne h$. The $q$ supporters of $h'$ have issued notarization support for a block other than $h$, so none of the correct ones can final-vote $h$. The validators that can final-vote $h$ are therefore the at most $N-q=f+p$ that do not support $h'$, together with at most $f$ Byzantine supporters, giving at most
\begin{equation}
(N-q)+f=2f+p=q-1
\end{equation}
final votes for $h$: one short of the $q$ a finalization certificate requires. The guard is what makes the two sets disjoint, and no rule about the order of a validator's votes is needed. For fast finalization, at least $N-p-f'$ correct validators first-voted for $h$, where $f'\le f$ is the actual Byzantine count. At most $p$ correct first votes remain for other hashes, so no other block can reach the second-look threshold $f+p+1$ or a size-$q$ notarization quorum. $\square$

The same counting gives a fact used twice below: a voting domain that holds notarization certificates for two different values can finalize neither of them, because each certificate denies the other value the $q$ final votes it would need.

\paragraph{Lemma 3.2 (Epoch closure does not skip account positions).} A decided epoch state never creates an empty account block. If no block survives at an unresolved slot $v$, the next epoch's clean election remains at $v$. If conflicting checkpoint-notarized blocks survive at $v$, a recovery block is at $v+1$ and names one surviving block at $v$ as parent.

\emph{Proof.} Account slot is determined only by the parent relation in an owner-authenticated block. Epoch closure changes the status of represented blocks but creates no synthetic account block. With no surviving block at $v$, the finalized parent remains at $v-1$, so the next valid child is again at $v$. With surviving notarized blocks at $v$, each is a real represented parent at that slot, so a child that chooses one of them is at $v+1$. $\square$

\paragraph{Lemma 3.3 (Epoch-election preservation of finalized account evidence).} Let $h$ be an account block finalized by a valid account-domain certificate. No valid epoch proposal can decide a finalized application state that replaces $h$ with a conflicting block at the same account position.

\emph{Proof.} For normal finalization, a size-$q$ final certificate contains at least $q-f=f+p+1$ correct signers. Correct signers reject an epoch proposal that omits or conflicts with the finalized position. Excluding all of them leaves at most
\begin{equation}
N-(f+p+1)=2f+p=q-1
\end{equation}
possible old-committee votes, too few for old-committee notarization of that epoch value. For fast finalization, a size-$q_{\mathrm{fast}}=N-p$ first-vote certificate contains at least $N-p-f=2f+p+1=q$ correct first-voters that protect the finalized position, again preventing a conflicting old-committee quorum. Checkpoint-finalized blocks are already part of predecessor finalized application state, which every valid successor state extends. $\square$

\paragraph{Theorem 3.4 (Irreversible finality with ancestral closure).} Once a block is finalized---by an account finalization certificate, by ancestral closure from a finalized descendant, or by unique inclusion in a decided epoch state---no later valid account or epoch decision can replace it with a conflicting block at the same slot.

\emph{Proof.} Account-domain finalization is protected by Lemma~3.1 in its voting domain and by Lemma~3.3 across epoch transitions. Ancestral closure finalizes only a single parent chain whose unresolved ancestors are already notarized; when the descendant finalizes, incompatible siblings are excluded from the finalized projection. Cross-epoch closure is permitted only after every earlier ancestor epoch has closed, so the descendant cannot retroactively choose a block that the earlier epoch omitted. Checkpoint finalization is part of the jointly decided state and successor states must extend its finalized projection. Induction from genesis therefore gives an irreversible finalized prefix for every account. $\square$

\subsubsection{Report preservation and reconciliation}
\paragraph{Lemma 3.5 (Historical-state soundness).} If a validator accepts a reconstructed report state for signed root $r_i$, then, except with negligible cryptographic failure probability, it has reconstructed the canonical certified state committed by $r_i$. Every vote-notarized or vote-finalized annotation it uses has a valid block and sufficient constituent signed votes for the corresponding local certificate.

\emph{Proof.} The reporter signature fixes the report context and target root $r_i$. Starting from an available source state, the requester applies the received reconstructive difference and accepts the candidate state $T_i$ only if $\StateHash(T_i)=r_i$. Collision resistance prevents two different canonical certified states from being accepted for the same root except with negligible probability. This root check authenticates the reconstructed state; it does not by itself establish the semantics of a certified-status annotation. Such an annotation is used only after the corresponding block and constituent signed votes satisfy the required $\NC$, $\FF$, or $\FC$ predicate and placement rules. $\square$

\paragraph{Lemma 3.6 (Common-state reconciliation).} Let $r_i$ be a correct reporter's signed epoch-$e$ state root. If correct replicas continue gossiping old epoch-$e$ blocks and signed votes and retain report history, then any correct requester eventually reconstructs $r_i$ whenever the report-availability execution reaches a common descendant state known both to the requester and to a correct replica retaining $r_i$.

\emph{Proof.} The requester advertises a locally available source root and target $r_i$. Before a common descendant exists, replicas that do not know both states need not answer. Once some correct replica knows a source root also available to the requester---possibly a later common descendant---and retains the historical state for $r_i$, it returns the reconstructive difference from that source to $r_i$. Eventual gossip delivers the reply and referenced objects, and the requester accepts only after the reconstructed target hashes to $r_i$. If old-epoch evidence quiesces, all correct replicas converge to one final live certified root, which can serve as the common source root for every retained correct report. $\square$

\paragraph{Lemma 3.7 (Late vote-notarization evidence survives reporting).} Suppose the signed account votes issued before correct reporters stopped epoch-$e$ voting are sufficient to construct an epoch-$e$ notarization certificate for block $h$. For every selected set of $N-f$ usable reports, either $h$ is certified-visible or at least one selected correct report contains residual first/notarization support for $h$.

\emph{Proof.} Let $V$ be $q$ distinct signers of the notarization support and $Q$ the $N-f$ selected reporters. Their identity sets overlap in at least
\begin{equation}
q+(N-f)-N=f+p+1
\end{equation}
validators. At most $f$ are Byzantine, so at least $p+1\ge1$ identities in the intersection are correct. Each such reporter either already summarized the locally constructible certificate in its certified state or records its own unsummarized signed support in residual evidence. Hence $h$ satisfies the report-derived inclusion rule in every valid selection. $\square$

\paragraph{Lemma 3.8 (A correct vote cannot be silently discarded while its lock is active).} Suppose a correct old-committee validator rejects an epoch proposal because its derived checkpoint ledger omits a block for which that validator issued an epoch-$e$ account vote. Under the gossip assumption, a later correct epoch leader can expose that block to the deterministic derivation by selecting the objecting validator's usable report.

\emph{Proof.} The validator gossips its signed vote and rejection context. Its own report commits either to a certified representation of the block or to residual signed support. Selecting that report makes the block certified-visible or gives $M_Q(h)\ge1$, so $\Include_Q(h)$ holds. The later deterministic state therefore includes the block unless a finalized incompatible branch justifies exclusion. $\square$

\paragraph{Lemma 3.9 (Timely account finalization is mandatory).} If every correct old-committee validator incorporated a valid account-finalization certificate for $h$ before signing its report, every valid epoch state finalizes $h$ and its eligible unresolved ancestor prefix.

\emph{Proof.} Any $N-f$ selected reports contain correct reporters whose certified states expose the finalized annotation and its evidence. The state derivation treats final-visible status as mandatory, and Theorem~3.4 prevents a conflicting checkpoint result. $\square$

\subsubsection{Epoch-election safety}
\paragraph{Lemma 3.10 (Cross-committee conflict excludes election).} If $O_e$ has $\NC_O(h)$ and $N_e$ has $\NC_N(h')$ for $h\ne h'$ in the same epoch-election slot, no epoch value can be elected in that slot.

\emph{Proof.} Electing a value $x$ requires $\FC(x)$ or $\FF(x)$ in both domains. By the counting of Lemma~3.1, a domain holding a notarization certificate for a value different from $x$ admits at most $q-1$ final votes for $x$, so it cannot produce $\FC(x)$; and $\FF(x)$ needs $N-p$ first votes for $x$, which leaves at most $p<f+p+1$ first votes for any other value, so that value would never reach the second-look threshold and could not have been notarized in that domain at all. Since $h\ne h'$, the value $x$ differs from $h$ or from $h'$, so at least one of the two domains can finalize nothing. Note that this covers a domain that has certified several values, which can finalize none of them, so the clause remains sound when one domain certifies a value the other has certified as well. The two notarization certificates are therefore safe skip evidence, and they are permanent: certificates are monotone, so a slot that the clause abandons never becomes decidable later. $\square$

\paragraph{Theorem 3.11 (Epoch-election agreement).} Two valid joint epoch-election finalization certificates for epoch $e$ cannot decide different states.

\emph{Proof.} Let epoch value $X$ be elected in epoch-election slot $w$. It is finalized in both epoch-election voting domains. Standard Kudzu safety in those timeout-bearing domains excludes a timeout certificate or conflicting notarization that could justify skipping the decided slot. A later epoch placement therefore descends from $X$. The fixed-value rule preserves $X$'s $(Q_e,d_e)$ along that lineage, so any later joint finalization decides the same report-derived state hash. Collision resistance and deterministic derivation identify the same $S_e$. $\square$

\subsubsection{Transfer safety}
\paragraph{Theorem 3.12 (Transfer safety).} Correct validators' decided finalized projections contain compatible account histories. Balances remain nonnegative, every send identifier $(h(B),j)$ is received at most once, and finalized account positions are never replaced.

\emph{Proof.} Induct over closed epochs from genesis. Epoch-election agreement gives one decided state for each predecessor. Only finalized blocks are replayed into $\Sigma_e$; checkpoint-notarized conflicting branches in $L_e$ have no application effects. Finalization and ancestral closure select one compatible parent chain, and every finalized block is owner-authenticated. Its ordered operation batch is valid only if the running balance never becomes negative. A receive names a concrete send identifier and is valid only if the send is addressed to the receiving account, is already finalized, and is absent from the finalized parent history and earlier operations of the same block. Theorem~3.4 prevents later replacement of any finalized position. $\square$

Domain separation prevents replay of old voting evidence. An epoch-$e$ account vote does not satisfy an epoch-$(e+1)$ account predicate. A block may be proposed again after a clean rollback, but it requires fresh votes. Equivocation recovery instead advances to a fresh child at the next slot, whose parent explicitly identifies the surviving notarized branch.

\subsubsection{Account progress without a scheduler}
RAI has no account leader schedule, round-robin service rule, slot timer, or bound on the number of other accounts that must be serviced first. An inactive account consumes no voting slot. The owner of an active account directly gossips its block whenever it chooses.

Consider a non-equivocating owner that gossips one valid block $B$ with hash $h$ during a sufficiently long synchronous portion of an open epoch. Once one correct validator receives $B$, gossip delivers the block and its first vote to all correct validators within the synchrony bound. Because there is no account timeout, an earlier unrelated vote cannot force a correct validator to abandon the slot. Correct validators first-vote the same valid hash, and the ordinary notarization path yields enough signed support to construct $\NC(h)$. If $B$ is finalization-eligible, correct validators can then issue final votes and construct $\FC(h)$. If it is not yet finalization-eligible only because an earlier parent epoch has not closed, it may remain notarized until that closure barrier is satisfied.

If an epoch boundary arrives before account finalization, the attempt is not converted into a timeout. The decided epoch state handles the unresolved slot: a unique compatible included block is checkpoint-finalized; multiple conflicts are checkpoint-notarized; an omitted provisional block is rolled back; and an empty unresolved position becomes a clean election in a later epoch.

\paragraph{Theorem 3.13 (Non-equivocating owner).} Let $B$ lie on an owner's valid chain. Suppose the owner eventually stops equivocating and continues gossiping the chain. Assume recurring sufficiently long synchronous portions of open epochs, eventual closure of required predecessor epochs, report availability, gossip, and epoch-election progress. Then every finite-depth block on the chosen chain is eventually finalized and checkpointed. No account scheduling or account-slot-timeout premise is required.

\emph{Proof.} Work forward from the latest finalized ancestor. In a clean slot, persistent gossip during a sufficiently long open interval makes the owner's next block vote-notarized and, when the ancestral closure conditions hold, account-finalized. If the epoch closes first, a unique surviving block is checkpoint-finalized. If the current parent is one of several checkpoint-notarized survivors from an earlier closed epoch, the owner advances to the next slot with a child naming that parent; finalization of the child recursively finalizes the chosen unresolved ancestor prefix. Each finite predecessor therefore eventually becomes finalized, after which later checkpoints preserve it by Theorem~3.4. $\square$

\paragraph{Theorem 3.14 (Recovery after finite equivocation).} Suppose an owner creates only finitely many conflicting blocks and eventually chooses one represented valid branch and stops equivocating. If a closed epoch leaves conflicting checkpoint-notarized blocks on that branch's unresolved frontier, the owner can recover by extending the chosen survivor at the next account slot. Under the report-availability, gossip, synchrony, and epoch-election-progress assumptions, the chosen branch is eventually finalized and incompatible provisional siblings are rolled back.

\emph{Proof.} Let $A_v$ be the chosen checkpoint-notarized frontier block in a closed epoch, with one or more incompatible checkpoint-notarized siblings at slot $v$. The owner creates a fresh block $A_{v+1}$ naming $A_v$ as parent. Because the earlier epoch is closed, $A_v$ satisfies the cross-epoch side of the ancestral-closure rule. During a sufficiently long later open epoch, $A_{v+1}$ becomes vote-notarized and then finalized, or is uniquely checkpoint-finalized. Either event finalizes $A_v$ and any unresolved notarized ancestors on the chosen branch and removes incompatible siblings. The argument repeats only if the chosen represented branch itself contains a deeper finite unresolved fork. $\square$

\subsubsection{Epoch-election progress}
At least $N-f$ old-committee validators are correct. If they reach the boundary, retain their signed report roots and committed residual objects, preserve enough authenticated state history to reconstruct those roots, and serve reconstruction, an epoch leader can obtain $N-f$ usable reports without Byzantine help. Common-state reconciliation may temporarily return no answer when no replica knows both requested roots. Under the report-availability premise, gossip eventually creates a common descendant root that the requester can use as a new source; a responder retaining that source and the historical target supplies the reconstructive difference, and the requester verifies success by recomputing the target root.

Once $S_{e-1}$ is known, $\BuildState$ deterministically derives a valid checkpoint ledger and finalized projection from any selected usable report set, subject to local preservation checks. A proposal rejected because it omits a correct validator's voted block exposes signed evidence. Including that validator's usable report in a later selection makes the block satisfy $\Include_Q$; the deterministic state then includes it provisionally, finalizes it if unique, or excludes it only if a finalized incompatible branch already resolves the slot. Because the closed account epoch contains finitely many signed account votes, only finitely many distinct preservation objections can arise.

\paragraph{Lemma 3.15 (Every synchronous epoch-election slot advances).} During a sufficiently long synchronous interval, an epoch-election slot eventually yields either a joint-complete epoch placement or $\ET$.

\emph{Proof.} Run the ordinary Kudzu progress argument separately in $O_e$ and $N_e$. In each timeout-bearing epoch-election domain, correct validators eventually receive enough signed votes to construct either a non-timeout notarization certificate or a timeout certificate. A timeout in either domain gives $\ET$. Otherwise both domains have non-timeout notarization certificates. If each domain has certified exactly one value and the two agree, the epoch placement is joint-complete and can still be decided; in every other case the two domains hold certificates for different values and the cross-committee conflict clause gives $\ET$. Preservation rejections expose finite report content that a later correct leader can select and deterministically incorporate. $\square$

\paragraph{Theorem 3.16 (Epoch-election liveness).} Assume $N-f$ correct reports become usable through terminating state reconciliation, recurring sufficiently long synchronous intervals, gossip of valid preservation-objection evidence, and an epoch-leader schedule with recurring correct leaders. Then every closing epoch eventually closes.

\emph{Proof.} The report phase eventually provides at least $N-f$ usable reports. If a correct validator rejects a report selection because the derived checkpoint ledger omits a locally protected voted block, it gossips the signed reason. A later correct leader can include the objecting validator's report, causing the block to satisfy $\Include_Q$. Since there are finitely many old-epoch votes, after finitely many such discoveries a correct leader can select reports whose deterministic state passes the preservation checks of enough correct old-committee validators. Lemma~3.15 then ensures each synchronous epoch-election slot either advances the common lineage or yields safe skip evidence. Eventually a correct leader proposes the valid report set and derived state hash; correct members of both committees reconstruct the same state, support the same epoch value, and locally construct matching finalization certificates. The epoch closes. $\square$

These progress guarantees are deliberately separate. Account elections themselves have no slot timers. Account progress needs sufficient gossip time before an epoch boundary; report reconstruction needs historical-state availability; cross-epoch account finalization waits for earlier unresolved epochs to close; and epoch closure needs epoch-leader progress. An owner that continues equivocating can keep its own account provisional.

\paragraph{Cost and limits.} RAI borrows Kudzu's account-vote quorum logic but removes per-account timeout voting. The joint epoch election still uses Kudzu timeouts. Parallel account elections remove a shared account-ordering dependency, while ancestral finalization lets one finalized descendant resolve a notarized prefix. Validators still replicate and verify blocks and votes. Historical reconciliation requires persistent authenticated state or deltas, repeated equivocation can enlarge the checkpoint ledger, and local preservation locks require old-committee validators to retain their own epoch vote evidence until the epoch election decides.

\section{Implementation and Evaluation}
\noindent\textbf{Reserved: two pages.} No RAI source code, experiment logs, or measurements accompanied the protocol specification. This section is an experiment plan and space allocation, not a report of completed work.

\subsection{RsNano-based prototype}
The intended prototype uses RsNano, a Rust implementation of a Nano node~\cite{RsNano}. The implementation description should identify the exact baseline revision, reused components, and changes for batched account blocks, timer-free account-slot voting domains, checkpoint-notarized fork state, ancestral finalization, child-based fork recovery, certified-state roots, residual vote evidence, historical state differences, local epoch-preservation locks, and joint epoch elections. It should state any differences in cryptography, storage, block format, admission, or networking.

The artifact should expose separate events for account-block receipt, vote-notarization, checkpoint-notarization, ancestral finalization, rollback of provisional forks, report signing, certified-state-root changes, state-difference replies, target-root reconstruction checks, preservation-lock rejections, and closed-state checkpointing. Persistent storage should distinguish the checkpoint ledger $L_e$, finalized application projection $\Sigma_e$, live certified trees, report-time historical roots, immutable report-time residual vote records, local old-epoch vote locks, and the deltas plus authentication hashes required to reconstruct historical state. A reconciliation reply should contain only the reconstructive difference material; it should not introduce a separately signed proof object. The account-election implementation must not create per-slot timers or account timeout votes. The epoch-proposal path should expose selected report roots and the proposed derived-state hash separately; it should not serialize the full $S_e$ as an epoch-proposal field.

\begin{center}
\fbox{\begin{minipage}{0.9\columnwidth}
\textbf{Reserved: implementation description.} Insert the verified architecture, baseline and prototype commit identifiers, changed modules, storage and recovery rules, and artifact instructions. No implementation-completion claim is made here.
\end{minipage}}
\end{center}

\paragraph{Correctness tests before benchmarking.} Test multiple sends and receives in one block, duplicate receive identifiers inside one block, double receive across finalized blocks, conflicting owner blocks at one account slot, long account-election delays without any account timer, unique epoch inclusion causing checkpoint finalization, conflicting epoch inclusion causing checkpoint notarization without application effects, omission and rollback of provisional blocks, clean re-election after an empty unresolved slot, same-epoch descendant finalization of unresolved ancestors, rejection of cross-epoch ancestral finalization before the earlier epoch closes, recovery after closure by creating a new block at slot $v+1$ that references one conflicting checkpoint-notarized block at slot $v$, old-vote replay, late local certificate assembly, state-root evolution after report signing, reconstruction through a later common state root, rejection of a malformed difference whose reconstructed state does not hash to the target root, preservation of immutable report-time residual objects, omission of already summarized votes from later residual views, and overlapping and disjoint epoch-election committees. The split-election regression should send different valid epoch proposals to the two committees and verify that validators which collect the underlying votes can locally construct conflicting notarization certificates and derive $\ET$.

\subsection{Experimental setup}
Use the same machines, resource limits, network topology, and offered operation traces for RAI and the RsNano baseline. Document processor, memory, storage, operating system, compiler, cryptographic choices, bandwidth, delay distribution, and peer fanout. Pin code revisions and configuration files. State warm-up length, run duration, independent repetitions, random seeds, and uncertainty reporting.

Vary active accounts, operations per block, send/receive mix, committee size, $f$, $p$, epoch duration, offered load, fork rate, the fraction of unresolved forks that require child-based recovery after an epoch boundary, and the depth of notarized ancestor prefixes finalized by one descendant. Control validator omission, equivocation, selective delivery, and offline periods. Runs presented as theorem-covered liveness experiments must satisfy the stated fault, synchrony, report-availability, and lock-objection-delivery premises; beyond-bound stress tests should be labeled separately.

\begin{table}[t]
\centering
\caption{Required comparison scenarios. All results are pending.}
\begin{tabularx}{\columnwidth}{@{}lX@{}}
\toprule
Scenario & Controlled change\\
\midrule
No forks & All owners follow one valid chain.\\
Forks & A chosen fraction of owners propose competing blocks at the same account slot.\\
Forks + Byzantine & Add validator equivocation, withheld votes, or selective delivery.\\
Forks + offline & Add unavailable validators and later recovery.\\
Ancestral finality & Allow a descendant to finalize several unresolved notarized ancestors; vary same-epoch depth and cross-epoch recovery depth.\\
\bottomrule
\end{tabularx}
\end{table}

Reconfiguration is an additional dimension. Vary committee overlap, deliberately delay historical-state reconstruction and earlier epoch elections, and test periods immediately before and after epoch boundaries. Distinguish owner equivocation from validator voting faults. There is no account scheduler parameter to tune; account load is generated by when owners inject blocks.

\begin{center}
\fbox{\begin{minipage}{0.9\columnwidth}
\textbf{Reserved: configuration table.} Insert system details, fixed and swept parameters, run counts, and workload definitions. Identify which runs satisfy the protocol assumptions.
\end{minipage}}
\end{center}

\paragraph{Comparable outcomes.} Report vote-notarization, checkpoint-notarization, irreversible RAI finalization, closed-state checkpointing, and the corresponding RsNano confirmation event separately. For batched blocks, report both blocks per second and operations per second, plus completed transfers per second. For end-to-end transfers, use the send identifier $(h(B),j)$ so a block containing several sends is not counted as one transfer. Report the latency from vote-notarization to finalization, the additional delay for cross-epoch closure barriers, and the number of unresolved ancestors finalized by each descendant.

\subsection{Latency and throughput results}
\noindent\textbf{Results pending.} Measure latency from owner block dissemination to vote-notarization, checkpoint-notarization when applicable, finalization, and checkpointing. For equivocation recovery, report the delay from closure of the forked epoch to finalization of the new child at slot $v+1$ and its selected ancestor at slot $v$. For clean rollback and re-election, report the number of epochs before a block at the unresolved slot finalizes. Report median, 95th percentile, and 99th percentile. For an end-to-end transfer, measure from send submission to the receive becoming finalized and checkpointed.

Plot useful throughput against offered load. Useful throughput counts unique operations or completed transfers at the stated persistence level; repeated clean-election attempts, losing provisional forks, and epoch-election timeout votes are overhead. Report backlog and unfinished operations instead of dropping them from latency statistics. Use identical workload seeds in paired RAI and baseline runs.

\begin{center}
\fbox{\begin{minipage}{0.9\columnwidth}
\textbf{Reserved: throughput and latency figure.} Insert measured load-throughput and latency curves for the required scenarios. Separate block throughput from operation throughput, show the cost of provisional-fork rollback and child-based recovery, and plot the number of unresolved ancestors finalized by a descendant.
\end{minipage}}
\end{center}

\paragraph{Resources.} Measure network bytes, signature verification work, processor utilization, and storage per finalized and checkpointed operation. Separate ordinary account traffic from report, state-difference, preservation-objection, and epoch-election traffic. Measure the size of certified-state history, authenticated deltas, residual vote evidence, checkpoint-notarized forks, and local epoch-preservation state. Compare the cost of a direct difference reply with the case where a reconciliation request initially has no responder and succeeds only after replicas reach a later common root.

\begin{center}
\fbox{\begin{minipage}{0.9\columnwidth}
\textbf{Reserved: resource results.} Insert bandwidth, computation, storage, reconciliation-history, and preservation-lock measurements with the same workload and completion definitions.
\end{minipage}}
\end{center}

\subsection{Faults, recovery, and reconfiguration}
\noindent\textbf{Results pending.} Measure progress for unaffected accounts while other owners fork. For recovering owners, report time from the last equivocation to closure of the forked epoch, creation and finalization of the recovery child at the next slot, recursive finalization of the selected ancestor prefix, and rollback of incompatible checkpoint-notarized siblings.

For every epoch election, measure report availability, time until target report roots become reconstructible, state-difference bytes, epoch-proposal validation, preservation-lock rejections, skipped epoch-election slots, and joint-election duration. Include low and zero committee overlap, delayed old committee members, offline recoveries, an epoch election that outlasts the next account epoch, deliberate split-epoch-leader executions, and delayed delivery of final votes sufficient to construct a finalization certificate locally, where the finalized parent already supports finalized descendants. Show how often reconciliation succeeds immediately from a replica that already knows both roots versus after convergence to a later common root.

\begin{center}
\fbox{\begin{minipage}{0.9\columnwidth}
\textbf{Reserved: fault and transition figure.} Insert account-progress, ancestral-finality, fork-recovery, report-reconstruction, and epoch-transition measurements. Show useful throughput and irreversible finality around reconfiguration, not only notarization throughput.
\end{minipage}}
\end{center}

\paragraph{Ablations and limitations.} Compare persistent authenticated-state differences with a full historical-state transfer fallback. Compare operation with and without committee changes while holding workload and hardware fixed. Compare strict checkpoint-gated finalization against RAI ancestral finalization to isolate the latency benefit of same-epoch closure and the recovery cost of the cross-epoch closure barrier. Sweep $p$, epoch duration, operations per block, and report-history retention policy. State any batching, cache, garbage-collection, or transport optimization that changes the base protocol.

\begin{center}
\fbox{\begin{minipage}{0.9\columnwidth}
\textbf{Reserved: interpretation and limitations.} Insert conclusions supported by measurements, adverse cases, and threats to comparison validity. Do not claim an improvement without a measured baseline.
\end{minipage}}
\end{center}

\section{Related Work}
\noindent\textbf{Reserved: approximately one page, shared with the conclusion.} The paragraphs below identify the directly used work. A complete comparison with asset-transfer protocols and Byzantine reconfiguration remains to be written.

\paragraph{Single-owner transfers.} Guerraoui et al.~\cite{ConsensusNumber} show that a single-owner asset-transfer object has consensus number one and give a Byzantine message-passing construction. This motivates separating independent account histories instead of imposing one total order on all transfers. RAI nevertheless uses elections to certify account-chain choices and to agree on reconfiguration checkpoints.

\paragraph{Account chains.} Nano's block-lattice design gives each account its own chain and separates sends from receives~\cite{Nano}. RAI follows this account-local organization but permits a block to batch several send and receive operations. RsNano is the intended implementation baseline~\cite{RsNano}, not evidence that the RAI mechanisms have already been implemented or measured.

\paragraph{Kudzu and gossip-based BFT.} RAI directly borrows Kudzu's quorum thresholds, first-vote structure, second-look tests, and normal and fast finalization paths~\cite{Kudzu}. Kudzu builds on the Simplex line~\cite{Simplex}. RAI deliberately removes timeout voting from account elections: an unresolved account domain simply remains unresolved until finalization or the epoch boundary. Kudzu timeout voting is retained only inside the joint epoch-election instances, where it provides safe skip evidence for failed reconfiguration slots. RAI also changes the ordered objects and account semantics by running account domains in parallel, allowing decided checkpoints to checkpoint-notarize conflicting forks without application effects, and using descendant finalization to resolve a notarized ancestor branch. Tendermint supplies the reference model for one gossip layer~\cite{Tendermint}, not RAI's voting logic.

\paragraph{Authenticated state reconciliation.} RAI's reports commit to certified ledger state rather than reproducing complete vote histories. Historical reconstruction is based on persistent authenticated state and a difference operation between known roots. A difference contains the edits and authentication hashes needed to reconstruct the target state; correctness is checked by recomputing the signed target root, so RAI does not require a separate reconciliation-proof object. Correct replicas continue to learn old-epoch signed votes after reporting and construct resulting certificates locally; later common roots act as source states from which signed historical roots can be reconstructed. Residual vote evidence contains reporter-local votes not yet summarized into certified state, while local epoch-preservation locks retain the reporter's own old-epoch voting obligations until the epoch decision. The decided epoch ledger then either finalizes a unique survivor, checkpoint-notarizes conflicts, or rolls back omitted provisional blocks. A complete related-work section should compare this mechanism with authenticated state synchronization, anti-entropy, fork-aware reconciliation, Byzantine reconfiguration schemes that transfer explicit vote logs, and protocols that distinguish local finality from later checkpoint inclusion.

\begin{center}
\fbox{\begin{minipage}{0.9\columnwidth}
\textbf{Reserved: complete related-work comparison.} Add primary-source comparisons for fork recovery, authenticated state reconciliation, Byzantine reconfiguration, account-local finality, and ancestral finality. State each protocol's finality event, fault model, and progress conditions.
\end{minipage}}
\end{center}

\section{Conclusion}
RAI uses immutable single-owner account blocks as account-election proposals. Account elections have no slot timers and no timeout votes. An owner may broadcast a valid block at any time while an epoch is open; validators first-vote, notarize, and final-vote when the relevant evidence becomes available. Epoch boundaries stop further account voting but do not create synthetic account blocks.

A decided epoch state separates the checkpoint ledger from its finalized application projection. If one compatible block survives at an unresolved slot, checkpoint inclusion finalizes it. If several conflicting blocks survive, they are checkpoint-notarized in the ledger but have no application effects. Omitted provisional blocks are rolled back, while finalized blocks are never removed. If no block survives, a later epoch begins a clean election at the same account slot.

Finalization has an ancestral closure rule. A finalized descendant can finalize unresolved notarized ancestors from the same epoch immediately. Across an epoch boundary, the earlier ancestor epoch must close first. This gives a direct recovery path after equivocation: if epoch $e$ closes with conflicting checkpoint-notarized blocks at slot $v$, the owner creates a new block at slot $v+1$ referencing one survivor. Finalization of that child selects and finalizes the chosen ancestor branch and rolls back incompatible provisional siblings.

Epoch reports remain compact authenticated state commitments with reporter-local residual signed vote evidence. Reconciliation uses common descendant roots as source states for reconstructing historical report states, and the receiver verifies each reconstruction by recomputing the signed target root; no separate reconciliation-proof object is needed. The joint epoch election, unlike account elections, retains Kudzu timeout voting; old and new committees vote in separate domains on the same report-derived state hash. Matching finalization certificates decide the checkpoint, while timeout or cross-committee conflict evidence safely skips failed epoch-election slots.

The design therefore separates three forms of progress: timer-free account voting, epoch-boundary resolution of provisional forks, and timeout-bearing joint reconfiguration. The remaining engineering questions are the storage and bandwidth cost of authenticated history, residual vote evidence, checkpoint-notarized forks, ancestral finalization, and child-based recovery under faults. Those claims require an implementation and measurements before submission.

\paragraph{Use of AI tools.} ChatGPT was used to draft and organize this paper from supplied notes and protocol specifications, assist with mathematical exposition and source checks, identify protocol edge cases, and prepare the LaTeX document. No experimental measurements were generated. The protocol choices, proofs, citations, and implementation claims require author review before submission.

\begin{center}
\fbox{\begin{minipage}{0.9\columnwidth}
\textbf{Reserved: final related-work and conclusion edits.} Use this remaining space to complete the comparison and replace draft-only reservations with claims established by the implementation and experiments. Retain an accurate AI-use disclosure.
\end{minipage}}
\end{center}

\begin{thebibliography}{9}
\bibitem{Tendermint}
Ethan Buchman, Jae Kwon, and Zarko Milosevic. 2019. \emph{The Latest Gossip on BFT Consensus}. arXiv:1807.04938v3. doi:10.48550/arXiv.1807.04938.

\bibitem{DLS}
Cynthia Dwork, Nancy Lynch, and Larry Stockmeyer. 1988. Consensus in the Presence of Partial Synchrony. \emph{J. ACM} 35, 2 (1988), 288--323. doi:10.1145/42282.42283.

\bibitem{ConsensusNumber}
Rachid Guerraoui, Petr Kuznetsov, Matteo Monti, Matej Pavlovic, and Dragos-Adrian Seredinschi. 2019. The Consensus Number of a Cryptocurrency. In \emph{Proceedings of the 2019 ACM Symposium on Principles of Distributed Computing}. ACM. doi:10.1145/3293611.3331589. Extended version: arXiv:1906.05574.

\bibitem{Nano}
Colin LeMahieu. n.d. \emph{Nano: A Feeless Distributed Cryptocurrency Network}. White paper. \url{https://content.nano.org/whitepaper/Nano_Whitepaper_en.pdf}. Undated edition; accessed September 22, 2026.

\bibitem{RsNano}
RsNano contributors. 2026. \emph{RsNano: A Rust Port of nano-node}. Software repository. \url{https://github.com/rsnano-node/rsnano-node}. Accessed September 22, 2026; experimental revision not yet specified.

\bibitem{Simplex}
Victor Shoup. 2024. Sing a Song of Simplex. In \emph{38th International Symposium on Distributed Computing (DISC 2024)}, LIPIcs Vol.~319, 37:1--37:22. doi:10.4230/LIPIcs.DISC.2024.37.

\bibitem{Kudzu}
Victor Shoup, Jakub Sliwinski, and Yann Vonlanthen. 2025. \emph{Kudzu: Fast and Simple High-Throughput BFT}. arXiv:2505.08771v2. doi:10.48550/arXiv.2505.08771. Version of September 29, 2025.
\end{thebibliography}

\end{document}
